\section{Capua: The Pastor (1602--1605)}

\subsection{Appointment and a one-day discrepancy}

In 1602 the metropolitan see of Capua fell vacant, and Clement~VIII
gave it to Bellarmine. Benedict~XVI's 2011 catechesis dates the
appointment 18~March 1602 (M, reported). The consecration was
performed by the pope himself---an honor all witnesses record---but
the day is transmitted in two forms that cannot both be right.
\work{Lux illa} says Clement consecrated him ``with his own hand''
on 20~April of that year (AAS~22, 595; M). The autobiography says
the pope consecrated him ``on the Second Sunday after Easter, when
the Gospel \latin{Ego sum pastor bonus} is read''
(D\"ollinger--Reusch, 41; A)---and in 1602 the Second Sunday after
Easter fell on 21~April, the date Benedict~XVI's catechesis also
gives. The autobiography's liturgical anchor (Good Shepherd Sunday,
a detail Bellarmine plainly cherished) is the better witness; the
decretal's ``20~April'' is recorded as the discrepancy it is. Two
days after the consecration he received the pallium, and the day
after that he left the palace, hid four days in the Roman College
to escape farewell visits, and departed for his church---a haste,
he notes with rare dryness, that ``astonished many, and the
Pontiff himself, since men of the Curia can scarcely be torn from
the Curia''; a fellow bishop consecrated the same day did not
leave Rome until the end of October (A).

\subsection{Three years' residence, measured}

He entered Capua on 1~May 1602, and by Ascension Day was in the
pulpit (A). The autobiography then does something rare in
episcopal memoirs of any age: it gives auditable numbers. In three
years of residence he made the visitation of the whole diocese
three times; held three diocesan synods; celebrated one provincial
council---the first in eighteen years; restored the cathedral and
the archbishop's palace at a cost of some thousands of gold pieces
in the first year; registered the poor families of the city and
sent them monthly allowances; and preached not only in the
appointed seasons but on feasts through the year, against the
local custom that confined cathedral preaching to Advent and Lent
(D\"ollinger--Reusch, 41--42; A). Paul~VI's 1967 letter making him
co-patron of the archdiocese repeats precisely this ledger---the
three synods, one provincial council, three visitations, the
catechesis, the defense of the poor---and quotes the
autobiography's own summary sentence: \latin{``amabatur a populo,
et ipse amabat populum''}, ``he was loved by the people, and he
loved the people'' (\work{Roma altera}, AAS~60, 12--13, citing
Autobiogr.\ c.~38; M quoting A). The royal ministers of Spanish
Naples, he adds, never troubled him, ``because they thought him a
servant of God'' (A).

The tenure was short but it was not nominal. Against the standing
suspicion---voiced in his own century and since---that the
appointment was a courteous exile from Rome after his \latin{de
auxiliis} counsels displeased Clement~VIII, the documents read
here are silent: no witness inspected for this study states
Clement's motive, and the question is preserved as open rather
than resolved by inference.

\subsection{The conclaves of 1605: ``From the Papacy deliver me''}

In March 1605 Clement~VIII died. Leo~XI reigned twenty-six days;
then came the conclave that elected Paul~V. ``In both conclaves,
especially that latter,'' says the 1907 Catholic Encyclopedia,
``the name of Bellarmine was much before the electors, greatly to
his own distress, but his quality as a Jesuit stood against him in
the judgment of many of the cardinals.'' The autobiography's
account of his own conduct is among its most quoted pages: he kept
to his cell or walked alone, praying, ``Send whom thou wilt
send,'' and ``\latin{A Papatu libera me Domine}''---``From the
Papacy deliver me, O Lord''; in the second conclave ``little was
lacking that he should be made Pope'' (\latin{parum abfuit, quin
fieret Papa}); when a person of great weight offered to work for
his election he urged him to desist, saying he would not lift a
straw from the ground to become pope; and he bore no grudge
against those who blocked him, ``for he said the papacy was a most
dangerous labor'' (D\"ollinger--Reusch, 42--43; A).

\witnesslimit{The near-election claim is Bellarmine's own
retrospective statement, seconded by the encyclopedia's report of
conclave talk; no conclave voting record was inspected for this
study, and D\"ollinger--Reusch's commentary itself weighs the
claim against the surviving conclave relations. It is reported
here at that ceiling: that he was seriously discussed is
multiply witnessed; how near he actually came is not
established.}

Paul~V kept the new-made elector in Rome. Bellarmine, ``obediently
complying,'' asked at least to be released from a see he could no
longer reside in, and resigned Capua (CE 1907); \work{Lux illa}
frames the same fact from the papal side---Paul~V ``retained him in
the Curia to use his wise counsels'' (AAS~22, 595--596; M). From
1605 until his death he lived in Rome, ``a member of the Holy
Office and of other congregations, and thenceforth the chief
advisor of the Holy See in the theological department of its
administration'' (CE 1907).
